%%%%%%%%%%  scan idx 206  |  folio 193  %%%%%%%%%%
\begin{equation}\tag{*}
 \mathbb R(Qf_*)f^*E\simeq \mathbb RQ\,f_*f^*E\simeq \mathbb R(f_*Q)f^*E.
\end{equation}
Comme $U$ est noethérien, $\Qcoh(U)$ possède assez d'objets injectifs dans $\Mod(U)$ ([H] II 7.18), donc le dérivé droit de $f_*:\Qcoh(U)\longrightarrow\Qcoh(S)$ est « induit » par $\mathbb Rf_*:\mathrm D^+(U)\longrightarrow\mathrm D^+(S)$ lorsqu'on identifie $\mathrm D^+(\Qcoh(U))$ (resp. $\mathrm D^+(\Qcoh(S))$) à une sous-catégorie pleine de $\mathrm D^+(U)$ (resp. $\mathrm D^+(S)$) grâce à (3.5). Par suite, on a
\[
 \mathbb R(f_*Q)f^*E\simeq \mathbb Rf_*\mathbb RQf^*E.
\]
On en déduit une suite spectrale
\[
 E_2^{pq}=R^pf_*R^qQf^*E\Longrightarrow R^{*}(f_*Q)f^*E.
\]
D'après le cas 1°), il existe un entier $n(U)$ ne dépendant que de $U$ tel que $Q$ soit de dimension cohomologique $\leq n(U)$ sur $\Mod(U)$. On a donc
\[
 E_2^{pq}=0\quad\text{pour }q>n(U).
\]
Mais on a
\[
 E_2^{pq}=0\quad\text{pour }p>N,
\]
à cause de l'hypothèse de « dimension cohomologique » sur $S$. On a donc
\[
 R^n(f_*Q)f^*E=0\quad\text{pour }n>N+N'',
\]
où $N''$ est un entier supérieur à tous les $n(U)$, $U$ parcourant l'ensemble des intersections $U_{i_0}\cap\ldots\cap U_{i_p}$. Tenant compte de (*), on obtient
\[
 R^nQ(f_*f^*E)=0\quad\text{pour }n>N+N'',
\]
d'où finalement
\[
 R^nQ(E)=0\quad\text{pour }n>N+N'+N'',
\]
